%!TEX root = ../main.tex
% APÉNDICE B — Detalles técnicos e implementación

\chapter{Detalles técnicos e implementación}\label{apx:tecnicos}

Este apéndice reúne la información necesaria para reproducir el sistema a
partir del repositorio (anexo~\ref{anx:repositorios}).

\section{Versiones de las tecnologías}

\begin{table}[H]
\centering
\caption{Versiones utilizadas}\label{tab:versiones}
\small
\begin{tabularx}{\textwidth}{l l >{\raggedright\arraybackslash}X}
\toprule
\tb{Componente} & \tb{Versión} & \tb{Uso} \\
\midrule
.NET & 10.0 & Plataforma del backend y de las pruebas \\
ASP.NET Core (API, Identity, JWT Bearer) & 10.0.12 & API REST, cuentas de usuario y autenticación por token \\
Entity Framework Core (SQL Server) & 10.0.12 & Acceso a datos y migraciones \\
Swashbuckle.AspNetCore & 10.2.3 & Documentación interactiva de la API \\
SQL Server Express & -- & Base de datos en el entorno de desarrollo \\
Azure SQL Database & -- & Base de datos prevista para la publicación (próximo paso) \\
React y React DOM & 19.2 & Interfaz de usuario \\
Vite & 8.3 & Servidor de desarrollo y empaquetado del frontend \\
oxlint & 1.81 & Análisis estático del frontend \\
xUnit & 2.9.3 & Pruebas automatizadas \\
Entity Framework Core (SQLite) & 10.0.12 & Base en memoria para las pruebas \\
\bottomrule
\end{tabularx}
\end{table}

\section{Esquema de datos}

El esquema se define con clases de C\# y se versiona con 18
\glspl{migracion} de Entity Framework Core. Las entidades principales son
las siguientes.

\begin{itemize}
    \item \texttt{ApplicationUser}: cuenta de usuario, con su plan y el nombre de su grupo.
    \item \texttt{PersonaAsistida}: perfil de la persona (tipo de discapacidad, obra social, localidad, necesidades) y su dueño.
    \item \texttt{PersonaCompartida}: integrantes del grupo que acompañan a cada persona, con su rol (cuidador u observador).
    \item \texttt{Medicamento}, \texttt{Turno}, \texttt{EntradaDiario} y \texttt{Documento}: registros del cuidado diario.
    \item \texttt{TramitePersona}: avance de cada paso de la guía de trámites.
    \item \texttt{RecursoRecomendable} y \texttt{ValoracionRecurso}: catálogo de recursos y valoraciones de los cuidadores.
    \item \texttt{ArticuloRecomendado}: artículos de especialistas por tipo de discapacidad.
    \item \texttt{ContactoEmergencia} y \texttt{Bienestar}: red de apoyo y autoevaluaciones del cuidador.
\end{itemize}

\section{Variables de configuración}

\begin{table}[H]
\centering
\caption{Variables de configuración del backend y del frontend}\label{tab:variables}
\small
\begin{tabularx}{\textwidth}{l >{\raggedright\arraybackslash}X}
\toprule
\tb{Variable} & \tb{Propósito} \\
\midrule
\texttt{ConnectionStrings:DefaultConnection} & Cadena de conexión a la base de datos \\
\texttt{Jwt:Key} & Clave para firmar los tokens. No se guarda en el repositorio \\
\texttt{Jwt:Issuer}, \texttt{Jwt:Audience} & Emisor y destinatario de los tokens \\
\texttt{Cors:OrigenesPermitidos} & Orígenes del frontend autorizados a llamar a la API \\
\texttt{Frontend:Url} & Dirección del frontend, usada en los correos \\
\texttt{Correo:Servidor} y relacionadas & Servidor de correo para recuperar la contraseña \\
\texttt{Almacenamiento:CarpetaDocumentos} & Carpeta donde se guardan los documentos subidos \\
\texttt{Demo:Cargar}, \texttt{Demo:Renovar} & Carga y actualización de los datos de demostración \\
\texttt{VITE\_API\_URL} & Dirección de la API que usa el frontend \\
\bottomrule
\end{tabularx}
\end{table}

\section{Ejecución local}

Requisitos: SDK de .NET 10, Node.js 22 o superior, SQL Server Express (instancia \texttt{.\textbackslash SQLEXPRESS}) y la herramienta de Entity Framework (\texttt{dotnet tool install --global dotnet-ef}). Antes del primer arranque hay que copiar \texttt{appsettings.Local.example.json} como \texttt{appsettings.Local.json} y completar una clave JWT de al menos 32 caracteres. Ese archivo está excluido del repositorio para que los secretos no se publiquen.

\begin{lstlisting}[caption={Pasos para levantar el sistema en un entorno local}, label={lst:ejecucion}]
git clone https://github.com/valenquiroz/Tesis.git
cd Tesis/backend/Cuida.Api
copy appsettings.Local.example.json appsettings.Local.json
(editar Jwt:Key en appsettings.Local.json)
dotnet ef database update
dotnet run --launch-profile http --Demo:Cargar=true
cd ../../frontend
npm install
npm run dev
\end{lstlisting}

El parámetro \texttt{--Demo:Cargar=true} carga las familias, la residencia, los recursos y los artículos de demostración. El README del repositorio detalla las cuentas de prueba y, en Windows, el archivo \texttt{iniciar.bat} levanta el backend y el frontend juntos.
